\section{Controlled tensors and the ordinary coefficient reduction}
\label{app:outer}
Write $\tau=t$ for the equal-field interchange exponent and
$\gamma=g=1/(2-t)$ for the native endpoint and aggregate rectangle
exponent proved above. Constants are fixed before the input length
varies. For native dimensions $M=\prod_i n_i$, set
$L=\lceil\log_2(2M)\rceil$ and $V=MD$. The generic permutation tariff
established below is
\begin{equation}
 O(VL^\gamma\polylog L).
\end{equation}
This appendix develops the outer construction of~\cite{release}; the
arithmetic and movement costs use the suppliers proved in this paper.
\subsection{A tensor perfect shuffle}

Write $n_i=2k_i$ with the first digit most significant. Define
\[
 S_i(k_i h_i+\ell_i)=2\ell_i+h_i,
 \quad h_i\in\{0,1\},\quad 0\le\ell_i<k_i.
\]
Also define $d_M(X)=2X\bmod(M-1)$ for $X<M-1$ and fix $M-1$.
Splitting the flat stream into its two halves and interleaving records
implements $d_M$ in $O(V)$ time.

\begin{lemma}[Post-shuffle endpoint correction]
\label{lem:shuffle}
The tensor $\bigotimes_i S_i$ and its inverse have cost
\[O(VL^\gamma\polylog L)\]
by Lemmas~\ref{lem:movement}
and~\ref{lem:endpoint}.
\end{lemma}
\begin{proof}
For $x_i=k_i h_i+\ell_i$ and incoming carry $c\in\{0,1\}$,
\[
 2x_i+c=n_i h_i+(2\ell_i+c).
\]
The outgoing carry is exactly $h_i$, independent of $c$.
End-around doubling therefore has output
$y_i=2\ell_i+h_{i+1}$, cyclically in $i$.
Consequently
\[
 \bigotimes_i S_i=r\circ d_M,\qquad
 r(y)_i=2\lfloor y_i/2\rfloor+(y_{i-1}\bmod2).
\]
This rotates low endpoint bits \emph{after} the flat shuffle; it does
not require a cheap rotation of the original high-half bits.

Index the axes from zero. Reverse endpoint bits on $[0,t-1)$ and then
on $[0,t)$. A bit at $i<t-1$ moves first to $t-2-i$ and then to $i+1$;
the last bit moves to zero. This is precisely $r$. For a reversal
$[a,b)$, a left-half interval $I=[c,d)$ pairs with
$J=[a+b-d,a+b-c)$, in reverse local order. Partition the left half
into such intervals of length $u\le h$, where
\[
 \ell_R=\lceil\log_2R\rceil,\qquad
 h=\min\{\lfloor t/16\rfloor,\lfloor L/(40\ell_R)\rfloor\}.
\]
The complement of two selected intervals has at most three linear
intervals. Since $t\ge16h$, one has at least $u$ axes; choose a
contiguous $u$-axis tile auxiliary $W$ there.

Place the two selected banks $X,Y$ and $W$ at the front in that order.
Extract them successively from the remaining axis list: at each step
exchange the whole preceding aggregate gap with the requested interval.
At most three adjacent aggregate rectangle exchanges suffice. All
unselected axes keep their relative order; reverse the same exchanges
after the batch. Each exchange crops its two aggregate pads before the
next one and costs $O(VL^\gamma)$. No endpoint gather is performed.

With $\pi(j)=u-1-j$, the chronological endpoint masks
$Y_{\pi(j)}\mathrel{\oplus}=X_j$,
$X_j\mathrel{\oplus}=Y_{\pi(j)}$,
$Y_{\pi(j)}\mathrel{\oplus}=X_j$
exchange the paired bits. Each reads only the opposite frozen bank
and uses $W$ as its separate dirty tile auxiliary; its weighted offset
uses the target bank's own native place values. Every quotient label
is fixed. The three banks use at most $3L/40$ address bits. After
padding the one remaining suffix to $2^s$ and borrowing a dock of
width $c\le L/40+1$, its low index supplies
$R_{\rm op}=2^{s-c}D\ge2^{9L/10-O(1)}D$ opaque bits per record.
This pays the weighted-mask metadata and exact boundary repair even
when $D=1$. There are no phase-role bits in these direct masks.
The completed masks restore the dock and fix the low suffix index,
so the suffix pad is cropped before the inverse layout.

Each mask costs $O(V(u\log R)^\tau+V)$.
Since $\gamma>\tau$ and $h\log R\le L/40$, a batch costs
$O(VL^\gamma)$. From $L-2<\log_2M\le t\log_2R$,
$h\ge L/(80\ell_R)$ for sufficiently large admissible $L$.
Together with $t\le L/(3\log_2L)$ this gives $O(\log R)$ batches.
Summing proves the bound. Reversing the actual word gives the inverse;
the bounded remaining range uses a direct finite branch.
\end{proof}

\subsection{Permutations on exact even dimensions}

An adjacent pair can select either $I$ or the exchange matrix $J$.
Over $\F_4$, both have a common word of exactly eight $H$ letters
interleaved with unit diagonals. Tensoring a slot uses the supplied
canonical kernel. The diagonal coefficient factors by local axes;
its lexicographic list is generated by prefix products in $O(M)$
fixed-field operations. The completed word is binary. Its lane pad
can be removed only at that completed boundary.

Odd-start adjacent matchings are conjugates by coordinatewise $+1$.
Explicitly, if $T(x)=x+1\bmod n$, an odd-start pair $(i,i+1)$ is
$TKT^{-1}$ with $K$ switching $(i-1,i)$. Set $K=I$ on every
irrelevant pair, in particular $(n-2,n-1)$ whose conjugate would
wrap to $(n-1,0)$. Thus the completed matching has no cross-boundary
or wrap switch, even when its temporary $K$ crosses a ragged parent
boundary.
At a synchronized odd-start stage, translate every axis in the fixed
target set, including axes whose matching is identity. Choose $K=I$
on those axes; conjugation then cancels their translations. Thus the
translated set, flat offset and source mask are independent of switch
settings, including any control address in the controlled compiler.
This selective translation costs $O(V)$: flat rotation by the sum of
the selected native place values agrees except when a selected digit
wraps. The exceptional fraction is at most $t/L^3\le L^{-2}$.
Rotate the good-source mask with the payload. For bad output slots,
compute $x=P^{-1}(y)$ and the correct destination $T(x)$, sort those
records by destination, and refill the holes. Since $P=T$ on good
sources and both are permutations, $P(B)=T(B)$. This does not assume
that the bad source set is invariant. Keys, record sorting and
schoolbook decoding cost
\[
 O(\delta ML(D+L)+\delta ML^2)=O(V).
\]
Here the $L^2$ exception-address bound does not perform $t$ full
$L^2$ divisions. If $b_i=\lceil\log_2n_i\rceil$, then
$\sum_i b_i\le2L$. Streaming a current at-most-$L$-bit numerator
through a $b_i$-bit remainder register divides by $n_i$ in
$O(Lb_i+L+b_i)$ tape steps. Summing gives $O(L^2)$ for all digits;
the weighted Horner reconstruction has the same bound.

The source-good mask also has a linear implementation. Starting with
one bit, append native factors in their physical order. For each bit
$a$ of the previous mask, print $a$ for the $n_i-1$ nonmaximal values
and zero at the maximal value if axis $i$ is selected; otherwise
print $n_i$ copies of $a$. Counter reset work is at most its ensuing
$n_i$-bit output block. Since the intermediate mask lengths grow by
factors at least two, their sum is below $2M$. Rotate that mask with
the flat data permutation, so it marks exactly the output holes $P(B)$.
Compute and copy a full key only at a bad record. A monotone output
counter otherwise advances amortized in $O(M)$ total time. Stable
binary radix sorting in $L$ passes costs
$O(KL(D+L))$ for $K\le\delta M$ bad records, including the full
payload and header scans. In each pass stage one fixed-width item,
remember its selected key bit, and copy the item to the corresponding
stable zero/one bucket. Concatenate the buckets and rewind; a constant
number of tapes and $O(K(D+L))$ steps suffice per pass, including
staging and returns. Binary counters have amortized linear cost over
each scan; their per-item resets cost at most that item's width.
Since the sorted destinations are exactly
$P(B)=T(B)$, refill bad holes in that order; there is no $L$-bit key
comparison at each good position. This establishes the stated bound
even for $D=1$. The negative translation uses nonzero-digit source
masks and the same proof.

\begin{lemma}[Ragged global-shuffle network]
\label{lem:ragged}
Every permutation of an original even $n$ has a word of $O(\log n)$
levels consisting of adjacent switches and the \emph{same} full $S_n$
and $S_n^{-1}$. No recursive odd-sized child is padded.
\end{lemma}
\begin{proof}
At a level, the residual permutation preserves \emph{each} interval
setwise in a public ordered partition into contiguous parents.
Pair positions from each parent's
start, leaving its last position unmatched if its size is odd.
Use the same output pairing. The bipartite multigraph has input
pair/singleton vertices, output pair/singleton vertices, and one edge
per packet. Even parents give even cycles. An odd parent gives one
path between its two singleton vertices, plus even cycles. That path
has an odd number of edges. Its endpoint edge colors can both be
forced to the global parity of the parent's last position.

Alternating edge colors give input switches $E_{\rm in}$ and output
switches $E_{\rm out}$ so that the residual has identity
\[
 \pi=E_{\rm out}S_n C S_n^{-1}E_{\rm in}.
\]
The full inverse shuffle stably splits global even and odd positions.
For original even $n$, write this split as
\[
 p_n(x)=S_n^{-1}(x)=
 \begin{cases}
 x/2,&x\text{ even},\\
 n/2+(x-1)/2,&x\text{ odd}.
 \end{cases}
\]
A parent interval $[a,b)$ therefore has the two image fragments
\[
 \left[\left\lceil a/2\right\rceil,\left\lceil b/2\right\rceil\right),
 \qquad
 \left[n/2+\left\lfloor a/2\right\rfloor,\,
       n/2+\left\lfloor b/2\right\rfloor\right).
\]
The even fragments concatenate to $[0,n/2)$ and the odd fragments
to $[n/2,n)$, skipping empty fragments.
The new contiguous block list is
\[
 [\text{all parents' even fragments}]
 [\text{all parents' odd fragments}].
\]
The endpoints depend only on $n$ and the parent partition, independently
of $\pi$ and its switch colors. The residual $C$ preserves each new
fragment setwise. The recursive call is on this single $n$-wire
residual with the whole updated partition. Thus each level adds only
a fixed number of slot types to one recursive word.
Track this order; do not gather siblings. Each fragment has size at most
the ceiling of half its parent. After $\lceil\log_2n\rceil$ levels all
blocks are singletons. Each adjacent switch stage splits into its
even-start and odd-start families. Short axes can continue passive
singleton levels; their remaining shuffles cancel on unwind.
\end{proof}

Synchronize the networks for all axes to $\lceil\log_2R\rceil$ levels.
Lemmas~\ref{lem:shuffle}--\ref{lem:ragged} and the adjacent-switch
words give a generic local-permutation tensor cost
\begin{equation}
 \label{eq:generic}
 O(VL^\gamma\polylog L).
\end{equation}
This retains the parameter of the supplied kernels: each adjacent-switch
word uses a fixed number of native-$H$ calls, each shuffle has the
same $\gamma$ tariff, and there are $O(\log R)$ levels. Translations
and fixed-field diagonal scans cost $O(V)$. Thus synchronization adds
only a fixed power of $\log L$. Public local tables and network
settings have polynomial-time sequential generation in $R,L$.

\subsection{Controlled tensors and the control gauge}

Partition axes into targets $T$ and controls $C$. For each control
address $c$ let $\pi_i(c)$ be a permutation on target axis $i$.
Controls are address coordinates, not input-dependent payload values.

\begin{lemma}[Controlled compiler]
\label{lem:controlled}
Let $G(R,L)\ge1$ bound the time to generate one outer $\F_4$ scalar
coefficient at a current record address, including address decoding,
control-gauge inversion, all required local tables and ragged settings.
Assume this is a fixed polynomial-time sequential procedure with fixed
tapes. Then the controlled map
\[
 (c,x_T)\longmapsto(c,(\pi_i(c)(x_i))_{i\in T})
\]
has bit cost
\[
 O\bigl((MDL^\gamma+MG(R,L))\polylog L\bigr)
\]
with fixed tapes and ordinary order. In particular $G(R,L)=O(D)$
gives \eqref{eq:generic}. For the explicit CRT families below,
$G(R,L)=O((R+L+2)^{12})$; thus $D\ge(R+L+2)^{64}$ suffices.
\end{lemma}
\begin{proof}
First construct the completed subset kernel
$K_T=H_T\otimes I_C$. Use a static common eight-$H$ word for $H$ versus
$I$, tensoring each inner letter across all axes. The control axes may
move \emph{inside} this implementation. At its completed boundary,
$K_T$ fixes them. Insert controlled scalar stages only between completed
$K_T$ calls. The outer eight-$H$ switch word is then valid on each
fixed-control fiber. Its coefficient is
$\kappa(x,c)=\prod_{i\in T}\kappa_i(x_i,c)\in\F_4^*$.
We do not assume factorization by control axes.

The fragment-endpoint formulas in Lemma~\ref{lem:ragged} make the
partition and shuffle topology independent of the control address.
Only the switch settings depend on $c$, so a single fixed slot
schedule applies to all fibers.
Use the full tensor shuffle on \emph{all} axes. Give $C$ identity
switches. At an outer slot their temporary physical value is
$q_e(c)=S_C^e(c)$, where $e$ is the signed number of completed forward
minus inverse shuffles. At each scalar stage enumerate the current
native lexicographic addresses and decode logical
$c=q_e^{-1}(x_C)$. Generate the target settings for this logical $c$.
At every odd-start stage use the same translation of all axes in the
fixed set $T$, with identity $K_i(c)$ on idle axes or fibers. The flat
approximation can temporarily carry into a control coordinate on bad
records. Complete its exact repair before reading any control-dependent
setting. Completed subset kernels and translations then preserve the
physical controls; a shuffle increments $e$ in its prescribed direction.
Induction over the word proves the fiber invariant. Descent and unwind
balance, so $e=0$ at exit and the controls are restored.

For the CRT affine families, one concrete generator decodes addresses
by schoolbook division, computes the control CRT sums, generates the
at-most-$R$ local table, and constructs the pair graph by sequential
scans. All axes, levels, records, and full-list lookups are bounded in
the following generator analysis by $O((R+L+2)^{12})$ time. Fix the dominating
polynomial $P_1=(R+L+2)^{64}$ before choosing $D$. For each record park
its $D$ bits locally, generate one coefficient, apply its fixed two-bit
map, and advance. In general this scalar stage costs $O(M(D+G))$;
visited work intervals and returns are included in the generator bound.
There are only $\polylog L$ outer stages, proving the displayed
general bound. The CRT generator just described has $G=O((R+L+2)^{12})$,
which the chosen word width dominates, so its scalar stage costs $O(V)$.
No whole-array scan is nested inside a per-record lookup. This also
aligns coefficients to the actual post-shuffle stream order.
\end{proof}

\begin{lemma}[Sequential control-coefficient generation]
\label{lem:coefficient-generator}
For the affine CRT permutations used below, a coefficient at one current
native address, at one specified network slot, can be generated on a fixed
number of tapes in $O((R+L+2)^{12})$ bit time, including decoding, work
initialization, erasure and head returns.
\end{lemma}
\begin{proof}
Let $T=R+L+2$. All integers representing a native address or a CRT
product have $O(L)$ bits. Schoolbook division, multiplication and
extended Euclidean arithmetic have fixed polynomial cost in that width.
Decode the current mixed-radix address, invert the signed shuffle gauge,
and reconstruct any control block $x_A$ by ordinary CRT. There are at
most $L$ axes, so repeated schoolbook operations and sequential scans
give the deliberately loose bound $O(T^6)$ for this step.

For each target axis, list the at most $R$ images of the affine
permutation $z\mapsto P_Az+x_A\bmod p$, fixing its dummy value.
The ragged network above constructs its settings from these explicit
lists. There are at most $R^2$ fragment and pair records even if every
recursion level is stored separately. A record contains $O(L)$ bits.
Determine each local switch by tracing the alternating pair graph,
always looking up an endpoint by a full sequential scan. Fewer than
$O(R^2)$ path steps per fragment and $O(R^2L)$ bits per lookup suffice.
Allowing a separate traversal for every fragment and every one of at
most $L$ axes costs $O(R^6L^2)$, which is $O(T^8)$. Endpoint arithmetic
and the fixed common eight-$H$ word then specify the requested scalar.
Multiply the at most $L$ nonzero $\F_4$ entries in its tensor list.

These are nested loops of the stated finite bounds. A list lookup,
scratch initialization, return and erasure each has a complete scan
already charged above; replacing any direct table access by such a scan
still fits $O(T^{12})$. The generator receives only the public address
and slot descriptor, and performs no scan of the payload array. The
bound is therefore a per-record bit cost $G$, giving $MG$ for $M$
records, rather than $MDG$ for $D$ payload bits per record.
\end{proof}

\subsection{Enlarging the cyclotomic word}

\begin{lemma}[Uniform host initialization]
\label{lem:host-initialization}
From the binary input length $n$, for all sufficiently large $n$ a
deterministic fixed-tape procedure constructs a $2$-admissible tuple
$(q_0,\ldots,q_e)$, its product $2n<N<4n$, the attached common-multiple
data, and $\alpha,\Phi_\alpha,\rho$ with $q_0\mid\alpha$ and
$\rho$ a principal $(N/q_0)$-root. Its bit time is $N^{o(1)}$,
including writing and verifying the factor and root data.
\end{lemma}
\begin{proof}
Apply \cite[Proposition 3.4]{host} at $p=2$ and target $2n$;
its tuple and auxiliary data are part of the constructed admissible
length. Apply Proposition 3.8 of that source to this tuple. Their
respective time bounds are
$2^{O((\log\log n)^2)}$ and $2^{O((\log\log N)^4)}$ at fixed $p=2$.
Both are $N^{o(1)}$ in the source's deterministic fixed-tape model.
Write $L_0=\lceil\log_2N\rceil$ and
$\ell=\lceil\log_2L_0\rceil$; its divisor has
$\log\alpha=O(\ell^4)$.

Verification uses only polynomial-size setup operations. Trial division
checks the emitted prime list and its product, since each prime is at
most $R_0=2^{\ell^2}$; monic polynomial arithmetic checks the generated
cyclotomic data. Put $K=N/q_0$. Check $\rho^K=1$ and
\[
 \gcd(\rho^{K/q_i}-1,\Phi_\alpha)=1\qquad(1\le i\le e).
\]
The squarefree quotient is a product of fields. These tests force
order $K$ in every component, hence principality by the geometric-sum
identity. Modular powers, gcds and all certificates cost a polynomial
in $\alpha,R_0,L_0$, which is $2^{O(\ell^4)}=N^{o(1)}$.
Schoolbook arithmetic suffices in this setup; it uses no call to the
proposed multiplier.
\end{proof}

If $n$ is implicit in end-marked coefficient lists, counting and copying
them costs $O(n)$ by amortized counters. The bounded initial length
range uses the direct branch; no factorization is supplied as advice.

Use the unconditional host of \cite{host}. Put
$L_0=\lceil\log_2N\rceil$, $\ell=\lceil\log_2L_0\rceil$ and
$R_0=2^{\ell^2}$. Its squarefree $N$ has prime factors in
$(L_0^3,R_0]$. The constructed divisor $\alpha$ uses $O(\ell^2)$
factors and satisfies
\[
 \log\alpha=O(\ell^4),\qquad
 -\log\frac{\varphi(\alpha)}\alpha=O(\ell^2/L_0^3).
\]
The distinguished prime $q_0$ divides $\alpha$. Proposition 3.8
constructs $\rho$, a principal $(N/q_0)$-root in
$\F_2[Y]/\Phi_\alpha(Y)$, as well as this cyclotomic polynomial.
Then $\omega=\rho^{\alpha/q_0}$ is a principal $m=N/\alpha$ root.
Indeed for $s=\alpha/q_0$ and $0<j<m$, the old $K=sm$ character sum
at exponent $sj$ is $s$ copies of the $m$-sum. Odd $s$ is one in
characteristic two, so that sum vanishes.

Append unused prime factors until $\alpha'=\alpha J$ satisfies
\[
 D'=\varphi(\alpha')\ge(R_0+L_0+5)^{64}.
\]
Only $O(\ell)$ factors are needed because each multiplies $D$ by more
than $L_0^2$. There are $\Omega(L_0/\ell^2)$ available factors.
The first-crossing choice implies
\[
 \log\alpha'=O(\ell^4),\qquad
 \varphi(\alpha')/\alpha'>1-1/L_0,
 \quad m'D'=\Theta(N),\quad m'=N/\alpha'.
\]
We do not claim the enlarged divisor meets the original definition's
literal numerical upper cutoff; these displayed estimates are sufficient.
We retain the stronger density estimate as well. The original divisor
has $O(\ell^2)$ prime factors, and enlargement adds only $O(\ell)$,
all greater than $L_0^3$. Thus
\begin{equation}
 \label{eq:strong-density}
 \theta:=1-\frac{D'}{\alpha'}
   =1-\prod_{p\mid\alpha'}(1-p^{-1})
   \le\sum_{p\mid\alpha'}p^{-1}
   =O(\ell^2/L_0^3).
\end{equation}
This follows from the actual divisor construction in
\cite[Proposition 3.8, proof]{host}, rather than just its
weaker admissible-divisor definition.

The substitution $Y\mapsto U^J$ defines a unital homomorphism
\[
 \F_2[Y]/\Phi_\alpha(Y)\longrightarrow
 S=\F_2[U]/\Phi_{\alpha'}(U),
\]
because $\Phi_{\alpha'}(U)$ divides $\Phi_\alpha(U^J)$.
Write $\iota$ for this homomorphism. It preserves principal-root
character sums. The required root is
$\omega'=\iota(\omega)^J$: it satisfies $(\omega')^{m'}=1$.
For $0<j<m'$, split the old $m=Jm'$ terms into $J$ equal periods
to obtain
\[
 0=\sum_{a=0}^{m-1}\iota(\omega)^{Jja}
   =J\sum_{b=0}^{m'-1}(\omega')^{jb}.
\]
Odd $J$ is one in characteristic two, proving principality.
Thus no new root oracle is needed.

Pad the remaining prime axes by $n_i=p_i+1$. Then
$M/m'=\prod_i(1+1/p_i)=1+O(L_0^{-2})$.
Since $\alpha'>L_0^{64}$, for sufficiently large $N$,
$L=\lceil\log_2(2M)\rceil\le L_0$ and $L=\Theta(L_0)$.
The native radix hypotheses and $D'\ge P_1(R_0+1,L)$ hold together.
Setup polynomial in $R_0,D',L_0$ is $N^{o(1)}$ and is charged once.

\subsection{The coupled CRT index permutation}

Write $m'=\prod_{i=1}^t p_i$, with least-significant-first digits
\[
 x=a_1+p_1a_2+\cdots+(p_1\cdots p_{t-1})a_t.
\]
Its physical lexicographic axis order is $[p_t]\cdots[p_1]$, already
ordinary integer index order. In one lexicographic scan embed these
records in the $[p_i+1]$ grid, preserving their order, and zero-fill
every address containing at least one dummy coordinate.

\begin{lemma}[Parallel CRT tree]
\label{lem:crt}
The permutation $a\mapsto(x\bmod p_i)_i$ and its inverse have cost
$O(m'D'L_0^\gamma\polylog L_0)$.
\end{lemma}
\begin{proof}
For a contiguous block $I=A\cup B$ with $A$ lower significance,
write $x_I=x_A+P_Ax_B$, $P_A=\prod_{i\in A}p_i$.
After converting both children to their local residues, the left
residues are already correct for $x_I$. Each right target becomes
\[
 z_i\longmapsto P_Az_i+x_A\pmod{p_i}.
\]
This is a permutation because $P_A$ is a unit modulo $p_i$.
The left residue address determines $x_A\in[0,P_A)$ by ordinary
integer CRT arithmetic. Fix the target dummy, and use the identity
when a relevant left control is dummy. This is an explicit bijection
on every fiber and preserves every dummy-pattern class at completion.

At each level of a balanced tree take the union of all right axes as
$T$, and all left/passive axes as $C$. The disjoint node merges are
one controlled tensor from Lemma~\ref{lem:controlled}. There are
$O(\log(t+1))$ levels. This schedule avoids two separate full-volume
child executions per node. Reverse the stages and affine maps for
the inverse.
\end{proof}

\subsection{Characteristic-two local Fourier compilation}

For an odd prime $p$, let $u\in S$ be a principal $p$-root. Every field
component has exact order $p$, so $1-u^j$ is a unit for $0<j<p$.
Set $H_j=\prod_{a=1}^j(1-u^a)$ and $H_0=1$. The Newton matrix
\[
 N_{j,k}=\prod_{a=0}^{k-1}(u^j-u^a),\quad j\ge k,
\]
is diagonal--lower-Toeplitz--diagonal, with
$N_{j,k}=H_jH_{j-k}^{-1}(-1)^k u^{k(k-1)/2}$.
The symmetric Newton factorization is
$F_p(u)=N\operatorname{diag}(N_{kk})^{-1}N^T$.
All pivots remain units over the finite etale algebra.

The explicit formal triadic circuit in
Lemma~\ref{lem:formal-convolution}, with one operand fixed, supplies a
linear convolution DAG of size $O(b\log^8b)$ and depth $O(\log^2b)$
over every commutative $\F_2$-algebra. Bounded-degree copy trees are
included. Appendix~\ref{app:local-words} also gives explicit constant/linear
gate rules for this formal specialization, with no Boolean
idempotence identity substituted in the coefficient ring.

The borrowed-coordinate mechanism of \cite{local} restores arbitrary
dirty workspace by replaying once with the true sources and once with
the sources suppressed. Its algebra works in characteristic two.
Inside a width-$v$ Toeplitz recursion, choose chunk size
$b=c v/\log^8v$, with $c$ a sufficiently small fixed constant. The
two chunks and all gate coordinates fit within the same $v$ positions.
There are $O(\log^{16}v)$ source/target chunk pairs, each using
$O(\log^2v)$ replay rounds. A balanced recursion gives
$O(\log^{19}p)$ rounds on exactly $p$ coordinates, as proved in
Lemma~\ref{lem:exact-width-toeplitz}. Its source, target and borrowed
ports are disjoint; output-copy wires are counted inside the DAG.
All borrowed values
return before the next chunk pair. The finite small range is direct.

Work in $B=S\otimes_{\F_2}\F_4$ and take central
$\eta^2+\eta+1=0$. Every coefficient $\mu\in B$ is the sum of two
units: let $e$ be one on the components where $\mu=1$ and zero
elsewhere, and put $d=1+e(\eta+1)$. Then $d$ and $\mu+d$ are units.
All component degrees equal $r_B=\operatorname{lcm}(\operatorname{ord}_{\alpha'}2,2)$;
thus $e=1+(\mu+1)^{2^{r_B}-1}$ can be prepared by polynomially many
small-algebra operations. For unit $t$,
\[
 U(t)=\begin{pmatrix}1&t\\0&1\end{pmatrix}
 =\operatorname{diag}(t,1)HJ\operatorname{diag}(t^{-1},1),
 \qquad U(\mu)=U(d)U(\mu+d).
\]
Here is one uniform slot word for the whole round. Pad each $p$ axis
by its one dummy and complete a partial shear matching by pairing
its unused coordinates with coefficient $\mu=0$. Order every used
pair as (target,source), and send these pairs to canonical adjacent
pairs by a local permutation. Set $u_0=d$, $u_1=\mu+d$ and
$D(t)=\operatorname{diag}(t,1)$. In rightmost-first matrix order,
\[
 U(\mu)=D(u_0)HJ D(u_0^{-1}u_1)HJ D(u_1^{-1}).
\]
All diagonal entries are units. An unused pair has $u_0=u_1=1$,
so its completed word is $(HJ)^2=I$; it need not be fixed at an
intermediate letter. Tensoring uses exactly two full canonical $H$
slots, two full $J$ slots and three axis-factorized diagonal slots.
The pair relabeling and its inverse each cost one generic tensor
permutation from \eqref{eq:generic}; charge the two $J$ slots there
as well. Thus a round has two native-$H$ calls and at most four
generic permutation calls. No separate selection of $H$ versus $I$
is needed on this local path. Each completed local axis map is
$F_{p_i}\oplus[1]$ on its native-valid range $n_i=p_i+1$.
The tensor preserves each per-axis dummy-pattern subspace; it need
not fix every full address containing a dummy, since other axes still
transform. Consequently an initially zero-filled dummy complement
stays zero at the completed tensor boundary. Keep the full padded
grid through these operations, including every intermediate letter.

\begin{lemma}[Independent runtime word arithmetic]
\label{lem:old-word}
There is a uniform fixed-tape algorithm for ordinary binary polynomial
products at length $v$ in $O(M_{\rm old}(v))$ time, with bound function
\[
 M_{\rm old}(v):=v\log^2(2v)\log\log(4v).
\]
It uses classical integer multiplication, independently of the proposed
main construction and of the formal convolution circuit supplier.
\end{lemma}
\begin{proof}
For $v\ge2$ set $k=\lceil\log_2(v+1)\rceil$ and lift
$\sum_{i<v}a_iX^i$ to $\sum_{i<v}a_i2^{ki}$, with $a_i\in\{0,1\}$.
Every integer convolution coefficient is at most $v<2^k$.
Consequently the product's $k$-bit chunks are the exact integer
coefficients, without carries between chunks. Their low bits give the
ordinary $\F_2$ product. The encoding, zero padding, chunk extraction
and optional bit-order reversals are sequential scans costing $O(vk)$.
Counter resets are charged to their ensuing chunks. The classical
integer algorithm of \cite{ss} takes
$O(B\log(2B)\log\log(4B))$ steps on a fixed multitape machine for
$B=vk$ bits, proving the bound. Marking visited work-tape intervals
allows head returns and cleanup within the same order of time.
Length one is direct. The bound function is stable under fixed-factor
length changes.
\end{proof}

Tensor scalar coefficients are prepared by prefix expansion in the
current native order, as detailed in Appendix~\ref{app:local-words}. The total
number of products is below $2M$, not one product per axis at every
record. A $B$
symbol has $2D'$ bits. Each $B$ multiplication or fixed-scalar
application costs $O(M_{\rm old}(2D'))=O(M_{\rm old}(D'))
=D'\polylog L_0$ bit steps using Lemma~\ref{lem:old-word} and the
explicit fixed-modulus reduction in Appendix~\ref{app:local-words}.
In contrast, the routing compiler's $\F_4$
scalars are fixed two-bit maps costing $O(D')$. These two costs must
not be conflated. Project back to $S$ only after the complete Fourier
tensor has returned to that subspace.
Each synchronized replay round has a fixed number of calls at the
generic/native $\gamma$ tariff and a fixed number of these scalar
stages. There are $O(\log^{19} R)=\operatorname{polylog}L_0$ rounds,
so the local Fourier tensor retains the same $\gamma$ exponent.

\subsection{Ordinary multiplication and the time ledger}

Use Lemma~\ref{lem:host-initialization} to construct
$N=\Theta(n)>2n-2$ and its certified factor list, then enlarge the
divisor as above. Initial input scans cost $O(n)$ and all host
initialization is included in the $N^{o(1)}$ setup term below.
The top Agarwal--Cooley isomorphism from \cite[Lemma 2.2]{host} is
\[
 X\longmapsto Y^cZ,\quad c=(m')^{-1}\bmod\alpha',
\]
from $\F_2[X]/(X^N-1)$ to
$\F_2[Y,Z]/(Y^{\alpha'}-1,Z^{m'}-1)$.
Its coprime factors $\alpha',m'$ both exceed two for sufficiently large
$n$: $\alpha'$ tends to infinity and $\log\alpha'=o(\log N)$.
It uses one classical rectangle transpose and cyclic row shifts in
$O(N\log\alpha')=N\polylog L_0$ bit time, in either direction.
Explicitly, an exponent $k=r+m'j$ goes to
$(Z\text{-index},Y\text{-index})=(r,cr+j\bmod\alpha')$;
this fixes the row-shift convention and its inverse.
Since $\alpha'$ is odd, $\Phi_{\alpha'}$ is coprime to
$\psi=(Y^{\alpha'}-1)/\Phi_{\alpha'}$.
Put $f=\Phi_{\alpha'}$ and $d=\deg\psi=\alpha'-D'$.
If $d=0$ the coefficient split is the identity. Otherwise use
Lemma~\ref{lem:fixed-monic-remainder} on each $\alpha'$-bit row
for both projections, including the possibly much shorter modulus
$\psi$. The latter's prepared reciprocal has length $\alpha'-d=D'$.
Choose fixed Bezout polynomials $uf+v\psi=1$, with
$\deg u<d$ and $\deg v<D'$. The idempotents
$e_f=v\psi$ and $e_\psi=uf$ have degree $<\alpha'$, and reconstruct
\[
 A=(A_f e_f+A_\psi e_\psi)\bmod(Y^{\alpha'}-1).
\]
Two padded length-$\alpha'$ word products and a linear cyclic fold
suffice for this reconstruction. Projections and reconstruction thus
  cost $O(M_{\rm old}(\alpha'))$ per row, and
  \[
    O(m'M_{\rm old}(\alpha'))=N\polylog L_0
  \]
  overall. All reciprocals,
Bezout coefficients and idempotents are public polynomial-$\alpha'$
setup, charged once. The two projected streams have total length
$m'(D'+d)=N$ bits.

For the large component let $\omega$ be the principal $m'$ root and
\[
 E_i=(m'/p_i)((m'/p_i)^{-1}\bmod p_i),\quad \omega_i=\omega^{E_i}.
\]
The CRT idempotents give
$\omega^{xy}=\prod_i\omega_i^{(x\bmod p_i)(y\bmod p_i)}$.
Lemma~\ref{lem:crt} and the local Fourier tensor therefore compute
the true ordinary cyclic DFT. Two forward transforms, point products
and one inverse compute the product in $S[Z]/(Z^{m'}-1)$.
Retain the full native grid through these completed operations;
its dummy complement remains zero at each such boundary and under
the point product, so it can be cropped after inverse routing.
The inverse uses inverse roots and reversed index routing; its odd
normalization is one in characteristic two.

Here the large coefficient ring obtained by projection is
$\F_2[Y]/\Phi_{\alpha'}(Y)$, identified with $S$ by renaming $Y$ to $U$.
This identification is distinct from the earlier setup homomorphism
$Y_{\rm old}\mapsto U^J$ from the smaller ring
$\F_2[Y_{\rm old}]/\Phi_\alpha(Y_{\rm old})$. That homomorphism constructs
the root; it is not applied a second time to the new input words.
In this section $\omega$ denotes the already constructed $\omega'$.
Indeed if $J>1$, $U^J$ has order $\alpha$, rather than $\alpha'$,
in every field component of $S$. Hence $\Phi_{\alpha'}(U^J)$ is a
unit, so that repeated substitution would not even descend from the
new quotient.

Frequency order needs its own verification. If local roots instead have
exponents $b_i h_i$, with $b_i=m'/p_i$ and $h_i$ a unit modulo $p_i$,
an output digit $s_i$ describes frequency residue $b_i h_i s_i$.
The idempotent choice $h_i=b_i^{-1}$ makes that residue $s_i$.
For a coordinate permutation $Q$, however, $A=QF$ still satisfies
$A^{-1}((Aa)\odot(Ab))=F^{-1}((Fa)\odot(Fb))$.
Thus matching inverse and product checks alone do not certify a
literal Fourier index convention.

The complement has $d=\deg\psi=\alpha'-D'$ and
$v=dm'=N\theta=O(N\ell^2/L_0^3)$ by \eqref{eq:strong-density}.
If $d=0$, omit this component.
Pack $Y^jZ^i\mapsto T^{2di+j}$ in $\F_2[T]$. There is no overlap
between coefficient-product slots because $j+j'<2d$.
Each input has degree at most $2v-d-1$ and the product degree at most
$4v-2d-2$. Use Lemma~\ref{lem:old-word} at input length $2v$.
Crop its known zero tail to $2d(2m'-1)$ coefficient slots and unpack
$2m'-1$ rows of $2d$ bits, each with final bit zero.
Reduce their degree-$\le2d-2$ polynomials modulo the fixed monic
$\psi$ using Lemma~\ref{lem:fixed-monic-remainder}, then fold
$Z^{m'}=1$. All packing, unpacking and folding scans cost $O(v)$;
when $d=1$ the row remainders are direct.
The total cost is
$O(M_{\rm old}(2v)+m'M_{\rm old}(d))
 =O(M_{\rm old}(v)+m'M_{\rm old}(d))$. Since
$d\le v\le N$, each logarithmic factor is bounded using $L_0$,
so the total is
\[
 O(vL_0^2\ell)=O(N\ell^3/L_0)=o(N).
\]
The scans and all reductions are included. The weaker estimate
$v<N/L_0$ alone would not pay for this integer-packing fallback;
the stronger density is essential here. No integer packing of the
full ordinary input is used.

\begin{lemma}[Typed top-level execution]
\label{lem:typed-execution}
The preceding component algorithms compose on a fixed number of streams,
with ordinary coefficient input and output. Their complete calls use
the stated tariffs; no coefficient projection or outer native-dummy crop
occurs inside a Fourier word.
\end{lemma}
\begin{proof}
Use the following order, with $a,b$ denoting the two inputs.
\begin{enumerate}
\item Count and copy the coefficient lists, construct the host and all
length-dependent tables, and zero-fill both lists to $N$ coefficients.
The live top-level input streams contain $2N$ bits. Their cost is
$O(n)+N^{o(1)}$ before the linear zero-fill.
\item Apply the top isomorphism to each list, producing $m'$ consecutive
$\alpha'$-bit coefficient rows. Project each row to the $f$ and $\psi$
components. The two large streams contain $m'$ words of $D'$ bits;
the two small streams contain $m'$ words of $d$ bits. Park the small
streams on a fixed separate pair of tapes. The maps and projections cost
$N\polylog L_0$. The large words are identified with $S$ by the variable
rename just specified.
\item Embed each large word in $B=S[\eta]/(\eta^2+\eta+1)$ as $a+0\eta$,
and zero-fill its native dummy grid. Keep these two $M$-word $B$ arrays
through both forward transforms, their point product and the inverse.
They occupy $4MD'=O(N)$ bits at these top-level boundaries.
Transform one while the other waits on its own tape; constant whole-array
passing and parking scans cost $O(N)$. No callee traverses the parked
array.
\item To specify every completed transform, let $E:B^{m'}\to B^M$
be the zero-fill embedding from integer mixed-radix order. Let $Q$ be
the completed extended CRT permutation, and put
\[
 \mathcal T_{\omega}=\bigotimes_i(F_{p_i}(\omega^{E_i})\oplus[1]),
 \qquad \mathcal D_{\omega}=Q^{-1}\mathcal T_{\omega}Q.
\]
The CRT exponent identity gives
\begin{equation}
 \label{eq:typed-dft}
 \mathcal D_{\omega}E=E F_{m'}(\omega).
\end{equation}
Thus the output is again in ordinary integer frequency order, in the
same native grid. The whole operators $Q$ and $\mathcal T_\omega$
preserve every dummy-pattern subspace. Their intermediate words may
mix the auxiliary coefficient lanes or dummy values; only completed
boundaries are used in this identity. Their valid block has entries
in $S$, so a complete transform returns $S$-valued words when its
input was $S$-valued. Keep the unused $\eta$ lane allocated nonetheless.
\item Replace the first completed transform array by the coordinatewise
$B$ product of the two arrays, scanning them together, and erase the
second array. On valid coordinates this is exactly the $S$ product;
on every dummy coordinate it is zero. Performing all $M$ products,
without an exceptional-address lookup, costs
$O(MM_{\rm old}(D'))=N\polylog L_0$. Apply
$\mathcal D_{\omega^{-1}}$ to the first array. The principal-root
character sums imply
$F_{m'}(\omega^{-1})F_{m'}(\omega)=m'I=I$.
Equation~\eqref{eq:typed-dft} and the convolution identity therefore
give the cyclic product in $S[Z]/(Z^{m'}-1)$.
Explicitly, for the projected coefficient vectors $a_f,b_f$ the
$j$th inverse coefficient is
\[
 \sum_{r,s=0}^{m'-1}(a_f)_r(b_f)_s
       \sum_{k=0}^{m'-1}\omega^{k(r+s-j)}
   =\sum_{r+s\equiv j\pmod {m'}}(a_f)_r(b_f)_s.
\]
The inner character sum is zero off the stated congruence and is
$m'=1$ on it. Thus the complete retained-grid computation is
$\mathcal D_{\omega^{-1}}
 ((\mathcal D_\omega E a_f)\odot(\mathcal D_\omega E b_f))
 =E(a_f*_{\rm cyc} b_f)$.
Only now remove the zero $\eta$ lane and crop the zero dummy complement,
using monotone native-address scans costing $O(N)$. The result is
$m'$ consecutive $D'$-bit words in ordinary $Z$ order.
Generate the dummy test once per record using the fixed degree-twelve
address procedure; its cost is paid by $D'\ge P_1$.
\item Compute the parked small-component product by the independent
packing and fixed-monic reductions above, at cost $o(N)$, obtaining
$m'$ consecutive $d$-bit words, omitting this branch if $d=0$.
Reconstruct the $\alpha'$-bit rows from
the two components and apply the inverse top isomorphism at cost
$N\polylog L_0$. Extract the first $2n-1$ coefficients in a final
$O(N)$ scan. All other top-level streams are erased and their heads
returned by scans charged to their lengths.
\end{enumerate}
There are three complete transforms, each with a constant number of CRT
maps and the stated synchronized local rounds. Their total is
$O(NL_0^\gamma\polylog L_0)$, since their $B$ payload width is $2D'$
and hence satisfies the same metadata lower bound. The number of
top-level streams is fixed; setup archives and recursive scratch use
the separately charged tape families, rather than one fresh tape per
axis, table or recursion level.
\end{proof}

The full-volume routing and Fourier arrays have $V=\Theta(N)$ and
address width $L=\Theta(L_0)$. The smaller complement arrays have
the separate budget above. The stage counts are fixed powers of $\log L_0$.
The complete bound furnished by the construction is
\[
 O\bigl(NL_0^\gamma(\log L_0)^C\bigr)
   +N\polylog L_0+N^{o(1)}.
\]
